\documentclass[tikz,border=12pt]{standalone}
\usepackage{tikz}
\usepackage[fontset=none]{ctex}
\setCJKmainfont{Noto Sans CJK SC}
\setCJKsansfont{Noto Sans CJK SC}
\usetikzlibrary{shapes.geometric, arrows.meta, positioning, fit, backgrounds, calc, shadows}

\definecolor{drawBlueFill}{HTML}{DAE8FC}
\definecolor{drawBlueLine}{HTML}{6C8EBF}
\definecolor{drawGreenFill}{HTML}{D5E8D4}
\definecolor{drawGreenLine}{HTML}{82B366}
\definecolor{drawOrangeFill}{HTML}{FFE6CC}
\definecolor{drawOrangeLine}{HTML}{D79B00}
\definecolor{drawPurpleFill}{HTML}{E1D5E7}
\definecolor{drawPurpleLine}{HTML}{9673A6}
\definecolor{drawRedFill}{HTML}{F8CECC}
\definecolor{drawRedLine}{HTML}{B85450}
\definecolor{drawGreyFill}{HTML}{F5F5F5}
\definecolor{drawGreyLine}{HTML}{666666}

\pgfdeclarelayer{background}
\pgfsetlayers{background,main}

\begin{document}
\begin{tikzpicture}[
    every node/.style={font=\footnotesize},
    base_box/.style={rectangle, rounded corners=3pt, align=center,
        minimum height=0.9cm, minimum width=2.5cm,
        drop shadow={opacity=0.15}, thick},
    blue_node/.style={base_box, fill=drawBlueFill, draw=drawBlueLine},
    green_node/.style={base_box, fill=drawGreenFill, draw=drawGreenLine},
    red_node/.style={base_box, fill=drawRedFill, draw=drawRedLine},
    grey_node/.style={base_box, fill=drawGreyFill, draw=drawGreyLine},
    purple_node/.style={base_box, fill=drawPurpleFill, draw=drawPurpleLine},
    arrow/.style={-{Stealth[scale=1.0]}, thick, color=black!70},
    arrow_data/.style={-{Stealth[scale=1.2]}, very thick, color=drawOrangeLine},
    arrow_dash/.style={-{Stealth[scale=1.0]}, dashed, thick, color=black!50},
    tag/.style={font=\scriptsize, fill=white, inner sep=2pt, rounded corners=1pt},
    title_box/.style={font=\bfseries\normalsize, align=center,
        minimum width=5.2cm, minimum height=0.7cm, rounded corners=4pt, thick},
]

% ===== 标题 =====
\node[title_box, fill=drawRedFill!35, draw=drawRedLine!60] at (2.8, 0.8) {传统集中式方案};
\node[title_box, fill=drawGreenFill!35, draw=drawGreenLine!60] at (13.5, 0.8) {本文联邦学习方案};

% ===== VS =====
\draw[drawGreyLine, thick, dashed] (8.0, 0.3) -- (8.0, -1.2);
\draw[drawGreyLine, thick, dashed] (8.0, -4.8) -- (8.0, -6.2);
\node[circle, fill=drawOrangeFill, draw=drawOrangeLine, minimum size=1.3cm,
      thick, font=\bfseries\large, text=drawOrangeLine,
      drop shadow={opacity=0.2}] at (8.0, -3.0) {VS};

% ======== 左侧 ========
\node[blue_node] (sA) at (0.5, -0.6) {数据源 A};
\node[blue_node] (sB) at (0.5, -2.0) {数据源 B};
\node[blue_node] (sC) at (0.5, -3.4) {数据源 C};

\node[red_node, minimum width=2.8cm, minimum height=2.4cm] (ctr) at (4.8, -2.0)
    {中心服务器\\{\scriptsize 汇聚全部数据}};

\draw[arrow_data] (sA.east) -- (ctr.west|-sA);
\draw[arrow_data] (sB.east) -- (ctr.west);
\draw[arrow_data] (sC.east) -- (ctr.west|-sC);
\node[tag, text=drawOrangeLine] at (2.6, -0.35) {原始数据};
\node[tag, text=drawOrangeLine] at (2.6, -1.75) {原始数据};
\node[tag, text=drawOrangeLine] at (2.6, -3.15) {原始数据};

\node[circle, fill=drawRedFill, draw=drawRedLine, minimum size=0.5cm,
      thick, font=\scriptsize\bfseries, text=white] at (4.8, -0.5) {!};

\node[grey_node] (mOld) at (4.8, -5.0) {训练模型};
\draw[arrow] (ctr) -- (mOld);

\node[font=\scriptsize, text=drawRedLine, align=left, anchor=north west] at (-0.2, -6.4) {%
    \textbullet{} 数据隐私泄露风险高\\
    \textbullet{} 单点故障\\
    \textbullet{} 传输带宽压力大\\
    \textbullet{} 合规风险};

% ======== 右侧 ========
% 节点 x=11.2, 宽2.5, 左边缘=9.95
\node[blue_node] (nA) at (11.2, -0.6) {本地节点 A};
\node[blue_node] (nB) at (11.2, -2.0) {本地节点 B};
\node[blue_node] (nC) at (11.2, -3.4) {本地节点 C};

\node[green_node, minimum width=2.8cm, minimum height=2.4cm] (agg) at (15.5, -2.0)
    {安全聚合服务\\{\scriptsize 仅接收梯度}};

\draw[arrow] (nA.east) -- (agg.west|-nA);
\draw[arrow] (nB.east) -- (agg.west);
\draw[arrow] (nC.east) -- (agg.west|-nC);
\node[tag, text=drawGreenLine] at (13.3, -0.35) {加密梯度};
\node[tag, text=drawGreenLine] at (13.3, -1.75) {加密梯度};
\node[tag, text=drawGreenLine] at (13.3, -3.15) {加密梯度};

\node[circle, fill=drawGreenFill, draw=drawGreenLine, minimum size=0.5cm, thick]
    (chk) at (15.5, -0.5) {};
\draw[drawGreenLine!80!black, very thick] (15.38, -0.53) -- (15.47, -0.63) -- (15.62, -0.4);

\node[purple_node] (mNew) at (15.5, -5.0) {全局模型};
\draw[arrow] (agg) -- (mNew);

% ===== 模型下发回路 =====
% 总线 x=9.2 (距节点左边缘 9.95 有 0.75cm 间距)
% 从 mNew 底部 → 底部横线 → x=9.2 竖线 → 分叉进节点
\def\busX{9.2}
\coordinate (fb_bot) at (15.5, -5.8);
\coordinate (fb_corner) at (\busX, -5.8);

% 主干：mNew底部 → 底部横线 → 向上竖线（带圆角转弯）
\draw[dashed, thick, color=black!50, rounded corners=6pt]
    (mNew.south) -- (fb_bot) -- (fb_corner) -- (\busX, -0.6);

% 三条分叉箭头进入节点
\draw[arrow_dash] (\busX, -0.6) -- (nA.west);
\draw[arrow_dash] (\busX, -2.0) -- (nB.west);
\draw[arrow_dash] (\busX, -3.4) -- (nC.west);

% 标注在底部横线上
\node[tag, text=black!55] at (12.4, -5.55) {模型下发};

% 优势列表
\node[font=\scriptsize, text=drawGreenLine, align=left, anchor=north west] at (10.2, -6.4) {%
    \textbullet{} 原始数据不出本地\\
    \textbullet{} 去中心化，无单点故障\\
    \textbullet{} 传输量小（仅梯度）\\
    \textbullet{} 天然满足数据合规};

% ===== 背景 =====
\begin{pgfonlayer}{background}
    \fill[drawRedFill!6, rounded corners=10pt] (-0.7, 1.3) rectangle (6.5, -8.0);
    \fill[drawGreenFill!6, rounded corners=10pt] (8.7, 1.3) rectangle (17.3, -8.0);
\end{pgfonlayer}

\end{tikzpicture}
\end{document}
